% ============================================================
% Configuração dos cabeçalhos e rodapés
% Teoria de Linguagens e Compiladores
% ============================================================

\usepackage{fancyhdr}

% Espaço suficiente para o cabeçalho
\setlength{\headheight}{14pt}

% ------------------------------------------------------------
% Estilo principal das páginas
% ------------------------------------------------------------

\pagestyle{fancy}

% Remove configurações padrão
\fancyhf{}

% Linha do cabeçalho
\renewcommand{\headrulewidth}{0.4pt}

% Linha do rodapé
\renewcommand{\footrulewidth}{0pt}

% ------------------------------------------------------------
% Páginas pares
% ------------------------------------------------------------

\fancyhead[LE]{\small\thepage}
\fancyhead[LO]{\small\leftmark}

% ------------------------------------------------------------
% Páginas ímpares
% ------------------------------------------------------------

\fancyhead[RE]{\small\leftmark}
\fancyhead[RO]{\small\thepage}

% ------------------------------------------------------------
% Evita que o título do capítulo seja convertido para
% letras maiúsculas automaticamente
% ------------------------------------------------------------

\renewcommand{\chaptermark}[1]{%
  \markboth{CAPÍTULO\ \thechapter.\ #1}{}
}

% ------------------------------------------------------------
% Estilo das páginas de abertura de capítulo
% ------------------------------------------------------------

\fancypagestyle{plain}{%
  \fancyhf{}
  \renewcommand{\headrulewidth}{0pt}
  \renewcommand{\footrulewidth}{0pt}
  \fancyfoot[C]{\thepage}
}

% =========================================================
% Ambientes didáticos - Teoria de Linguagens e Compiladores
% =========================================================

\usepackage[most]{tcolorbox}
\usepackage{xcolor}

% =========================================================
% Caixas para os ambientes matemáticos do Quarto
% =========================================================

% ---------------------------------------------------------
% Definição
% ---------------------------------------------------------

\tcolorboxenvironment{definition}{
  enhanced,
  breakable,
  colback=blue!2,
  colframe=blue!50!black,
  boxrule=0.7pt,
  arc=1.5mm,
  left=5mm,
  right=5mm,
  top=3mm,
  bottom=3mm,
  before skip=10pt,
  after skip=10pt
}

% ---------------------------------------------------------
% Teorema
% ---------------------------------------------------------

\tcolorboxenvironment{theorem}{
  enhanced,
  breakable,
  colback=violet!2,
  colframe=violet!50!black,
  boxrule=0.7pt,
  arc=1.5mm,
  left=5mm,
  right=5mm,
  top=3mm,
  bottom=3mm,
  before skip=10pt,
  after skip=10pt
}

% ---------------------------------------------------------
% Lema
% ---------------------------------------------------------

\tcolorboxenvironment{lemma}{
  enhanced,
  breakable,
  colback=green!2,
  colframe=green!45!black,
  boxrule=0.7pt,
  arc=1.5mm,
  left=5mm,
  right=5mm,
  top=3mm,
  bottom=3mm,
  before skip=10pt,
  after skip=10pt
}

% ---------------------------------------------------------
% Proposição
% ---------------------------------------------------------

\tcolorboxenvironment{proposition}{
  enhanced,
  breakable,
  colback=orange!2,
  colframe=orange!55!black,
  boxrule=0.7pt,
  arc=1.5mm,
  left=5mm,
  right=5mm,
  top=3mm,
  bottom=3mm,
  before skip=10pt,
  after skip=10pt
}

% ---------------------------------------------------------
% Corolário
% ---------------------------------------------------------

\tcolorboxenvironment{corollary}{
  enhanced,
  breakable,
  colback=cyan!2,
  colframe=cyan!45!black,
  boxrule=0.7pt,
  arc=1.5mm,
  left=5mm,
  right=5mm,
  top=3mm,
  bottom=3mm,
  before skip=10pt,
  after skip=10pt
}

% ---------------------------------------------------------
% Exemplo
% ---------------------------------------------------------

\tcolorboxenvironment{example}{
  enhanced,
  breakable,
  colback=gray!3,
  colframe=gray!55!black,
  boxrule=0.7pt,
  arc=1.5mm,
  left=5mm,
  right=5mm,
  top=3mm,
  bottom=3mm,
  before skip=10pt,
  after skip=10pt
}